\section{总结与延伸}

Week 5 通过 Warp CEO 的视角，从产品设计的高度审视了 AI 开发者工具的发展方向。七大产品原则不仅适用于终端工具，也是评估和选择任何 AI 开发工具的通用框架。

\subsection{关键要点}

\begin{itemize}
    \item 成功的 AI 开发工具从开发者已知的界面出发，而非要求学习全新范式
    \item ``5 minutes to WOW''——零上手摩擦是关键
    \item Chat/自然语言正在从辅助功能变为核心交互方式
    \item MCP 是工具生态的标准化基础
    \item 快速反馈循环和可解释性是不可妥协的产品要求
    \item Agent 自主工作流正在成为主流，但 human-in-the-loop 仍然重要
    \item Agent 配置标准将从碎片化走向统一
\end{itemize}

\subsection{课程中的产品判断表}

\begin{table}[H]
\centering
\begin{tabular}{p{0.24\textwidth} p{0.3\textwidth} p{0.34\textwidth}}
\toprule
判断问题 & 课程给出的方向 & 对实践者的含义 \\
\midrule
入口界面是什么 & 从熟悉界面出发 & 不要为了 AI 牺牲已有工作流 \\
核心交互是什么 & 从编辑走向对话与委托 & 让用户表达意图，而非机械描述实现 \\
标准会不会统一 & 会逐步统一到规则文件与工具协议 & 早点建设可共享的 agent 配置与项目约束 \\
价值如何体现 & 以更快反馈和更少摩擦体现 & 第一体验要能在几分钟内证明价值 \\
风险在哪里 & 自主执行与错误决策 & 需要权限控制、确认门槛和审计链路 \\
\bottomrule
\end{tabular}
\caption{Week 5 的产品判断压缩表}
\end{table}

\subsection{拓展阅读}

\begin{itemize}
    \item Warp 官网：\url{https://www.warp.dev}
    \item AGENTS.md 提案：\url{https://github.com/anthropics/agents-md}
    \item Cursor：\url{https://cursor.sh}
    \item Bolt.new：\url{https://bolt.new}
\end{itemize}

%% --- Fin del contenido --- %%

\end{document}
